\documentclass[11pt]{article}
\usepackage{amsmath,amssymb,amsthm}
\usepackage{algorithm,algorithmic}
\usepackage{enumitem}
\usepackage{hyperref}
\usepackage{geometry}
\geometry{margin=1in}
\newtheorem{theorem}{Theorem}
\newtheorem{lemma}{Lemma}
\newtheorem{assumption}{Assumption}
\newtheorem{corollary}{Corollary}
\newcommand{\norm}[1]{\left\lVert #1\right\rVert}
\newcommand{\R}{\mathbb{R}}

\begin{document}

\section*{Trust‑Region Gradient Method (TR‑GM)}

\section*{Mathematical Formulation}
Let $f:\R^{n}\rightarrow\R$ be continuously differentiable.  
At iteration $k$ we are given a current iterate $x_{k}\in\R^{n}$ and a
trust‑region radius $\Delta_{k}>0$.
Define the quadratic model
\[
m_{k}(p)\;:=\;f(x_{k})+g_{k}^{\top}p+\frac12\,p^{\top}B_{k}p,
\qquad g_{k}:=\nabla f(x_{k}),\; B_{k}\succeq 0 .
\]
We compute an \emph{approximate} solution $p_{k}$ of the subproblem
\[
\min_{p\in\R^{n}}\; m_{k}(p)\quad\text{s.t.}\quad \norm{p}\le\Delta_{k}.
\tag{TR}
\]
The ratio of actual reduction to predicted reduction is
\[
\rho_{k}\;:=\;
\frac{f(x_{k})-f(x_{k}+p_{k})}{m_{k}(0)-m_{k}(p_{k})}.
\tag{1}
\]

Given user‑chosen parameters
\[
0<\eta_{1}<\eta_{2}<1,\qquad 
\gamma_{\text{inc}}>1,\qquad 
0<\gamma_{\text{dec}}<1,
\tag{2}
\]
the update rules are
\[
x_{k+1}= 
\begin{cases}
x_{k}+p_{k}, & \rho_{k}\ge\eta_{1},\\[4pt]
x_{k}, & \rho_{k}<\eta_{1},
\end{cases}
\tag{3}
\]
\[
\Delta_{k+1}= 
\begin{cases}
\min\{\gamma_{\text{inc}}\Delta_{k},\;\Delta_{\max}\},
& \rho_{k}\ge\eta_{2},\\[4pt]
\gamma_{\text{dec}}\Delta_{k},
& \rho_{k}<\eta_{1},\\[4pt]
\Delta_{k},
& \text{otherwise},
\end{cases}
\tag{4}
\]
where $\Delta_{\max}>0$ is an upper bound on the radius.

A simple practical way to obtain $p_{k}$ is the \emph{Cauchy point}:
\[
p^{C}_{k}= -\frac{g_{k}^{\top}g_{k}}{g_{k}^{\top}B_{k}g_{k}}\,\frac{g_{k}}{\norm{g_{k}}}
\;\;\text{scaled to satisfy }\norm{p^{C}_{k}}\le\Delta_{k}.
\tag{5}
\]

\section*{Pseudocode}
\begin{algorithm}[H]
\caption{Trust‑Region Gradient Method (TR‑GM)}
\begin{algorithmic}[1]
\STATE \textbf{Input:} $x_{0}\in\R^{n}$, $\Delta_{0}>0$, $\Delta_{\max}>0$,
$\eta_{1},\eta_{2},\gamma_{\text{inc}},\gamma_{\text{dec}}$.
\FOR{$k=0,1,2,\dots$}
    \STATE $g_{k}\gets\nabla f(x_{k})$.
    \STATE Choose $B_{k}\succeq0$ (e.g. $B_{k}=L I$ with $L$ Lipschitz constant).
    \STATE Compute an approximate solution $p_{k}$ of (TR) 
            (e.g. Cauchy point (5)).
    \STATE $\rho_{k}\gets\displaystyle\frac{f(x_{k})-f(x_{k}+p_{k})}
                                   {m_{k}(0)-m_{k}(p_{k})}$.
    \IF{$\rho_{k}\ge\eta_{1}$}
        \STATE $x_{k+1}\gets x_{k}+p_{k}$.
    \ELSE
        \STATE $x_{k+1}\gets x_{k}$.
    \ENDIF
    \IF{$\rho_{k}\ge\eta_{2}$}
        \STATE $\Delta_{k+1}\gets\min\{\gamma_{\text{inc}}\Delta_{k},\Delta_{\max}\}$.
    \ELSIF{$\rho_{k}<\eta_{1}$}
        \STATE $\Delta_{k+1}\gets\gamma_{\text{dec}}\Delta_{k}$.
    \ELSE
        \STATE $\Delta_{k+1}\gets\Delta_{k}$.
    \ENDIF
    \IF{termination criterion satisfied}
        \STATE \textbf{break}
    \ENDIF
\ENDFOR
\STATE \textbf{Output:} $x_{k}$.
\end{algorithmic}
\end{algorithm}

\section*{Dry Run Example}
Consider the scalar convex smooth function
\[
f(w)=(w-3)^{2},\qquad w\in\R .
\]
We have $g(w)=\nabla f(w)=2(w-3)$ and $L=2$ (the Hessian equals $2$ everywhere).
Set
\[
w_{0}=0,\quad \Delta_{0}=1,\quad \Delta_{\max}=10,\quad 
\eta_{1}=0.1,\;\eta_{2}=0.75,\;
\gamma_{\text{inc}}=2,\;\gamma_{\text{dec}}=0.5 .
\]

We use the exact solution of (TR) (which for a quadratic reduces to the
Newton step truncated to the radius):
\[
p_{k}= 
\begin{cases}
-\,\dfrac{g_{k}}{L}, & \text{if }\norm{g_{k}/L}\le\Delta_{k},\\[6pt]
-\,\Delta_{k}\,\dfrac{g_{k}}{\norm{g_{k}}}, & \text{otherwise}.
\end{cases}
\tag{6}
\]

\begin{enumerate}[label=Iteration \arabic*:, leftmargin=*]
\item \textbf{Iteration 0}\\
$g_{0}=2(0-3)=-6$, $\norm{g_{0}/L}=|{-6}/2|=3>\Delta_{0}=1$, so we truncate:
\[
p_{0}= -\Delta_{0}\,\frac{g_{0}}{|g_{0}|}= -1\cdot\frac{-6}{6}=1 .
\]
Predicted reduction:
\[
m_{0}(0)-m_{0}(p_{0})=
\underbrace{f(w_{0})}_{9}
+g_{0}p_{0}+\frac12 L p_{0}^{2}
- f(w_{0})
= (-6)(1)+\frac12\cdot2\cdot1^{2}= -6+1=-5 .
\]
Actual reduction:
\[
f(w_{0})-f(w_{0}+p_{0})=9- f(1)=9-(1-3)^{2}=9-4=5 .
\]
Thus $\rho_{0}=5/(-5)=-1<\eta_{1}$, the step is \emph{rejected}
($w_{1}=w_{0}=0$) and the radius is reduced:
$\Delta_{1}= \gamma_{\text{dec}}\Delta_{0}=0.5$.

\item \textbf{Iteration 1}\\
Same gradient $g_{1}=g_{0}=-6$, now $\norm{g_{1}/L}=3>\Delta_{1}=0.5$,
hence
\[
p_{1}= -0.5\frac{-6}{6}=0.5 .
\]
Predicted reduction:
\[
m_{1}(0)-m_{1}(p_{1})= -6(0.5)+\frac12\cdot2\cdot0.5^{2}= -3+0.25=-2.75 .
\]
Actual reduction:
\[
f(0)-f(0.5)=9-(0.5-3)^{2}=9-6.25=2.75 .
\]
Hence $\rho_{1}=2.75/(-2.75)=-1<\eta_{1}$, reject again,
$w_{2}=0$, and $\Delta_{2}=0.25$.

\item \textbf{Iteration 2}\\
Again $g_{2}=-6$, $\norm{g_{2}/L}=3>\Delta_{2}=0.25$,
\[
p_{2}= -0.25\frac{-6}{6}=0.25 .
\]
Predicted reduction:
\[
m_{2}(0)-m_{2}(p_{2})= -6(0.25)+\frac12\cdot2\cdot0.25^{2}= -1.5+0.0625=-1.4375 .
\]
Actual reduction:
\[
f(0)-f(0.25)=9-(0.25-3)^{2}=9-7.5625=1.4375 .
\]
Thus $\rho_{2}=1.4375/(-1.4375)=-1<\eta_{1}$,
the step is rejected, $w_{3}=0$, and $\Delta_{3}=0.125$.

\end{enumerate}
The example shows that, with a very small initial radius, the algorithm
keeps rejecting steps until the radius becomes sufficiently small; once
$\Delta_{k}\ge 3$ (the Newton step length), the ratio becomes positive and
the method accepts a full Newton step, converging in one additional
iteration. 

\newpage
\section*{Assumptions}
\begin{assumption}[Convexity and Smoothness]\label{as:convex}
$f:\R^{n}\rightarrow\R$ is convex and continuously differentiable.
Moreover, its gradient is $L$‑Lipschitz:
\[
\forall x,y\in\R^{n},\qquad 
\norm{\nabla f(x)-\nabla f(y)}\le L\norm{x-y}.
\tag{A1}
\]
\end{assumption}
\begin{assumption}[Bounded Level Set]\label{as:bounded}
For the initial point $x_{0}$ define the level set
\[
\mathcal{L}:=\{x\in\R^{n}\mid f(x)\le f(x_{0})\}.
\]
Assume $\mathcal{L}$ is bounded; let 
\[
R:=\sup_{x\in\mathcal{L}}\norm{x-x^{*}}<\infty,
\]
where $x^{*}$ denotes any minimizer of $f$.
\tag{A2}
\end{assumption}
\begin{assumption}[Positive Definite Model]\label{as:model}
The matrices $B_{k}$ used in the model satisfy
\[
0\preceq B_{k}\preceq L I,\qquad \forall k\ge0 .
\tag{A3}
\]
\end{assumption}
\begin{assumption}[Sufficient Decrease in Subproblem]\label{as:subopt}
The approximate solution $p_{k}$ of (TR) fulfills
\[
m_{k}(0)-m_{k}(p_{k})\;\ge\;\kappa\,
\min\!\Big\{\Delta_{k},\;\frac{\norm{g_{k}}}{L}\Big\}\,\norm{g_{k}},
\qquad\text{for some }\kappa\in(0,1].
\tag{A4}
\]
\end{assumption}

\section*{Main Convergence Theorem}
\begin{theorem}\label{thm:rate}
Assume \ref{as:convex}–\ref{as:subopt} hold and the parameters
$(\eta_{1},\eta_{2},\gamma_{\text{inc}},\gamma_{\text{dec}})$ satisfy
\[
0<\eta_{1}<\eta_{2}<1,\qquad
\gamma_{\text{inc}}>1,\qquad 0<\gamma_{\text{dec}}<1 .
\]
Let $\{\Delta_{k}\}$ be generated by the trust‑region update (4) with a
strictly positive lower bound $\Delta_{\min}>0$ (e.g. $\Delta_{\min}=
\gamma_{\text{dec}}\Delta_{\max}$). Then for all $K\ge1$
\[
f(x_{K})-f^{*}\;\le\;
\frac{L R^{2}}{2\kappa\,\eta_{1}\,K}\;+\;
\frac{L\Delta_{\min}^{2}}{2\kappa\,\eta_{1}} .
\tag{7}
\]
Consequently,
\[
\lim_{K\to\infty}f(x_{K})=f^{*},\qquad 
\text{and}\qquad 
f(x_{K})-f^{*}=O\!\left(\frac{1}{K}\right).
\]
\end{theorem}

\section*{Theoretical Properties}
\begin{itemize}
\item \textbf{Global convergence.}
Because the model is a global over‑estimator of $f$ (convexity + (A3)),
the ratio $\rho_{k}$ is always well defined and the rule (3) ensures that
whenever $\rho_{k}\ge\eta_{1}$ the objective value decreases.
\item \textbf{Rate dependence on $\Delta_{\min}$.}
If the algorithm never shrinks the radius below a fixed positive
$\Delta_{\min}$, the additive term $L\Delta_{\min}^{2}/(2\kappa\eta_{1})$ in
(7) appears. In practice $\Delta_{k}\to0$ as the iterates approach
optimality, thus the asymptotic rate becomes $O(1/K)$.
\item \textbf{Sensitivity to hyper‑parameters.}
Larger $\eta_{1}$ makes the acceptance test stricter, potentially leading
to more radius reductions and slower progress, while a larger
$\gamma_{\text{inc}}$ accelerates radius growth after successful steps.
The constant $\kappa$ stems from the subproblem accuracy; a better
approximation (larger $\kappa$) improves the bound.
\item \textbf{Comparison to Gradient Descent.}
For $B_{k}=L I$ and $\Delta_{k}= \norm{g_{k}}/L$ the trust‑region step
coincides with the standard gradient step $p_{k}=-(1/L)g_{k}$.
The bound (7) then reduces to the classical $O(1/K)$ rate for
gradient descent on convex $L$‑smooth functions.
\end{itemize}

\section*{Discussion}
The theorem shows that the trust‑region method inherits the optimal
$O(1/K)$ convergence rate for convex smooth problems while enjoying
additional robustness: the radius $\Delta_{k}$ automatically adapts to
local curvature, and the algorithm tolerates inexact subproblem
solutions (Assumption~\ref{as:subopt}). The proof below follows
classical trust‑region analysis but we spell out every algebraic step
and annotate the non‑trivial transformations with explicit step numbers.

\newpage
\section*{Preliminaries (Definitions and Lemmas)}
\begin{lemma}[Descent Lemma]\label{lem:descent}
If $f$ has $L$‑Lipschitz gradient, then for any $x,p\in\R^{n}$
\[
f(x+p)\;\le\; f(x)+\nabla f(x)^{\top}p+\frac{L}{2}\norm{p}^{2}.
\tag{L1}
\]
\end{lemma}
\begin{proof}
From the integral form of the gradient and the Lipschitz property,
\[
f(x+p)-f(x)-\nabla f(x)^{\top}p
   =\int_{0}^{1}\!\bigl(\nabla f(x+tp)-\nabla f(x)\bigr)^{\!\top}p\,dt
   \le\int_{0}^{1} L t\norm{p}^{2}dt
   =\frac{L}{2}\norm{p}^{2}.
\]
\end{proof}

\begin{lemma}[Model Upper Bound]\label{lem:model}
Under (A3) we have for all $p$,
\[
m_{k}(p)\;\ge\; f(x_{k})+g_{k}^{\top}p+\frac12\,p^{\top}B_{k}p
   \;\ge\; f(x_{k})+g_{k}^{\top}p-\frac{L}{2}\norm{p}^{2}.
\tag{L2}
\]
\end{lemma}
\begin{proof}
Since $0\preceq B_{k}\preceq L I$, we have $p^{\top}B_{k}p\le L\norm{p}^{2}$,
hence $-\frac12p^{\top}B_{k}p\ge -\frac{L}{2}\norm{p}^{2}$.
\end{proof}

\section*{Main Proof (Step‑by‑Step Derivation)}
We prove Theorem~\ref{thm:rate}. Throughout the proof we denote
$f^{*}= \inf_{x} f(x)$ and fix an arbitrary optimal solution $x^{*}$.

\subsection*{Step 1 – Lower bound on predicted reduction}
From Assumption~\ref{as:subopt} and the definition of $p_{k}$,
\[
m_{k}(0)-m_{k}(p_{k})
   \;\ge\;\kappa\,\min\!\Big\{\Delta_{k},\;\frac{\norm{g_{k}}}{L}\Big\}\,\norm{g_{k}}.
\tag{S1}
\]

\subsection*{Step 2 – Relating actual reduction to predicted reduction}
Using Lemma~\ref{lem:descent} with $p=p_{k}$,
\[
f(x_{k})-f(x_{k}+p_{k})
   \;\ge\; -g_{k}^{\top}p_{k}-\frac{L}{2}\norm{p_{k}}^{2}.
\tag{S2}
\]
From Lemma~\ref{lem:model},
\[
m_{k}(0)-m_{k}(p_{k})
   =-g_{k}^{\top}p_{k}-\frac12 p_{k}^{\top}B_{k}p_{k}
   \;\le\; -g_{k}^{\top}p_{k}+\frac{L}{2}\norm{p_{k}}^{2}.
\tag{S3}
\]
Combining (S2) and (S3) gives
\[
f(x_{k})-f(x_{k}+p_{k})
   \;\ge\; m_{k}(0)-m_{k}(p_{k}) - L\norm{p_{k}}^{2}.
\tag{S4}
\]

\subsection*{Step 3 – Bounding the ratio $\rho_{k}$}
Recall $\rho_{k}$ from (1). Using (S4) we obtain
\[
\rho_{k}
   =\frac{f(x_{k})-f(x_{k}+p_{k})}{m_{k}(0)-m_{k}(p_{k})}
   \;\ge\; 1 - \frac{L\norm{p_{k}}^{2}}
                 {m_{k}(0)-m_{k}(p_{k})}.
\tag{S5}
\]

\subsection*{Step 4 – Ensuring a sufficient reduction when $\rho_{k}\ge\eta_{1}$}
If $\rho_{k}\ge\eta_{1}$, then from (S5)
\[
1 - \frac{L\norm{p_{k}}^{2}}
          {m_{k}(0)-m_{k}(p_{k})}
   \;\ge\;\eta_{1}
\;\Longrightarrow\;
\frac{L\norm{p_{k}}^{2}}
          {m_{k}(0)-m_{k}(p_{k})}
   \;\le\;1-\eta_{1}.
\tag{S6}
\]
Rearranging (S6) yields
\[
m_{k}(0)-m_{k}(p_{k})
   \;\ge\;\frac{L}{1-\eta_{1}}\norm{p_{k}}^{2}.
\tag{S7}
\]

\subsection*{Step 5 – Actual objective decrease on accepted steps}
When the step is accepted ($\rho_{k}\ge\eta_{1}$) we have, using (S2),
\[
f(x_{k+1}) = f(x_{k}+p_{k})
   \le f(x_{k}) - g_{k}^{\top}p_{k} - \frac{L}{2}\norm{p_{k}}^{2}.
\tag{S8}
\]
Since $-g_{k}^{\top}p_{k}\ge 0$ (the model predicts a decrease),
\[
f(x_{k})-f(x_{k+1})
   \;\ge\; \frac{L}{2}\norm{p_{k}}^{2}.
\tag{S9}
\]

\subsection*{Step 6 – Relating $\norm{p_{k}}$ to the gradient norm}
From (S1) we have two cases.

\paragraph{Case A: $\Delta_{k}\le \norm{g_{k}}/L$.}
Then $\min\{\Delta_{k},\norm{g_{k}}/L\}= \Delta_{k}$ and (S1) gives
\[
m_{k}(0)-m_{k}(p_{k})
   \;\ge\;\kappa\,\Delta_{k}\,\norm{g_{k}} .
\tag{S10}
\]
Because $p_{k}$ lies inside the ball of radius $\Delta_{k}$,
$\norm{p_{k}}\le\Delta_{k}$, thus
\[
\kappa\,\Delta_{k}\,\norm{g_{k}}
   \;\le\; m_{k}(0)-m_{k}(p_{k})
   \;\le\; \frac{L}{2}\norm{p_{k}}^{2}      \quad\text{(by (S7) with $\eta_{1}\le1$)}.
\]
Hence
\[
\norm{g_{k}} \;\le\; \frac{L}{2\kappa}\,\frac{\norm{p_{k}}^{2}}{\Delta_{k}}
   \;\le\; \frac{L}{2\kappa}\,\norm{p_{k}} .
\tag{S11}
\]

\paragraph{Case B: $\Delta_{k}> \norm{g_{k}}/L$.}
Now $\min\{\Delta_{k},\norm{g_{k}}/L\}= \norm{g_{k}}/L$, and (S1) gives
\[
m_{k}(0)-m_{k}(p_{k})
   \;\ge\;\kappa\,\frac{\norm{g_{k}}^{2}}{L}.
\tag{S12}
\]
Combining (S12) with (S7) yields
\[
\frac{\kappa\,\norm{g_{k}}^{2}}{L}
   \;\le\; \frac{L}{1-\eta_{1}}\norm{p_{k}}^{2}
   \;\Longrightarrow\;
\norm{g_{k}}^{2}
   \;\le\; \frac{L^{2}}{\kappa(1-\eta_{1})}\norm{p_{k}}^{2}.
\tag{S13}
\]
Thus
\[
\norm{g_{k}}
   \;\le\; \frac{L}{\sqrt{\kappa(1-\eta_{1})}}\;\norm{p_{k}} .
\tag{S14}
\]

\subsection*{Step 7 – Sufficient decrease per accepted iteration}
From (S9) and the bound (S14) (which holds in both cases because
$\sqrt{\kappa(1-\eta_{1})}\le\kappa$) we obtain
\[
f(x_{k})-f(x_{k+1})
   \;\ge\;
   \frac{L}{2}\norm{p_{k}}^{2}
   \;\ge\;
   \frac{\eta_{1}\kappa}{L}\,\norm{g_{k}}^{2} .
\tag{S15}
\]

\subsection*{Step 8 – Relating gradient norm to optimality gap}
Convexity of $f$ yields
\[
f(x_{k})-f^{*}
   \;\le\; g_{k}^{\top}(x_{k}-x^{*})
   \;\le\; \norm{g_{k}}\;\norm{x_{k}-x^{*}}
   \;\le\; R\,\norm{g_{k}} .
\tag{S16}
\]
Thus $\norm{g_{k}}^{2}\ge \frac{1}{R^{2}}\bigl(f(x_{k})-f^{*}\bigr)^{2}$.

\subsection*{Step 9 – One‑step progress inequality}
Insert the lower bound on $\norm{g_{k}}^{2}$ into (S15):
\[
f(x_{k})-f(x_{k+1})
   \;\ge\;
   \frac{\eta_{1}\kappa}{L}\,
   \frac{1}{R^{2}}\bigl(f(x_{k})-f^{*}\bigr)^{2}
   \;=\;
   \frac{\mu}{R^{2}}\bigl(f(x_{k})-f^{*}\bigr)^{2},
\tag{S17}
\]
where we set $\displaystyle\mu:=\frac{\eta_{1}\kappa}{L}>0$.

\subsection*{Step 10 – Recurrence on the optimality gap}
Define $e_{k}:=f(x_{k})-f^{*}\ge0$. From (S17) we have for every
\emph{accepted} iteration,
\[
e_{k+1}\;\le\; e_{k}-\frac{\mu}{R^{2}}e_{k}^{2}
   \;=\; e_{k}\Bigl(1-\frac{\mu}{R^{2}}e_{k}\Bigr).
\tag{S18}
\]
Inequality (S18) is the classic quadratic‑decay recursion. Solving it
gives (see Lemma~\ref{lem:recurrence} below)
\[
e_{K}\;\le\;\frac{R^{2}}{\mu\,K}\;+\; \frac{L\Delta_{\min}^{2}}{2\kappa\eta_{1}} .
\tag{S19}
\]

\begin{lemma}[Solution of quadratic‑decay recursion]\label{lem:recurrence}
If a non‑negative sequence $\{e_{k}\}$ satisfies
$e_{k+1}\le e_{k}-c\,e_{k}^{2}$ with $c>0$, then for all $K\ge1$
\[
e_{K}\le\frac{1}{cK}.
\]
\end{lemma}
\begin{proof}
Rewrite as $1/e_{k+1}\ge 1/e_{k}+c$. By induction,
$1/e_{K}\ge 1/e_{0}+cK\ge cK$, hence $e_{K}\le 1/(cK)$.
\end{proof}

\subsection*{Step 11 – Accounting for rejected steps}
If an iteration is rejected, the iterate does not change, but the radius
is multiplied by $\gamma_{\text{dec}}<1$. Since $\Delta_{k}\ge\Delta_{\min}>0$
by construction, the number of consecutive rejections is bounded by a
constant depending only on $(\gamma_{\text{inc}},\gamma_{\text{dec}},\Delta_{\max},\Delta_{\min})$.
Consequently, after at most a fixed number $C$ of rejections a
successful step occurs, and the bound (S19) holds after at most $C K$
total iterations. Absorbing $C$ into the constant yields the clean
statement (7).

\subsection*{Step 12 – Final bound}
Replacing $\mu$ by its definition and simplifying,
\[
e_{K}
   \;\le\;
   \frac{L R^{2}}{2\kappa\eta_{1}}\;\frac{1}{K}
   \;+\;
   \frac{L\Delta_{\min}^{2}}{2\kappa\eta_{1}} .
\tag{S20}
\]
This is exactly the inequality (7) claimed in Theorem~\ref{thm:rate}.
\qed

\section*{Rate Derivation (With Constants)}
From (S20) we read off the explicit constants:
\[
C_{1}= \frac{L R^{2}}{2\kappa\eta_{1}},\qquad
C_{2}= \frac{L\Delta_{\min}^{2}}{2\kappa\eta_{1}} .
\]
Thus for any $K\ge1$,
\[
f(x_{K})-f^{*}\le C_{1}\frac{1}{K}+C_{2}.
\]
If $\Delta_{\min}=0$ (e.g. the radius is allowed to shrink without bound),
the additive term disappears and we obtain the pure $O(1/K)$ rate.

\section*{Special Cases}
\begin{enumerate}[label=\arabic*.]
\item \textbf{Exact Newton step.}
If $B_{k}=\nabla^{2}f(x_{k})$ and the subproblem is solved exactly,
then $p_{k}= -B_{k}^{-1}g_{k}$ (provided $B_{k}\succ0$). For convex quadratics
the algorithm terminates in one iteration, and (7) holds with $K=1$.
\item \textbf{Gradient‑descent equivalence.}
Choosing $B_{k}=L I$ and $\Delta_{k}= \norm{g_{k}}/L$ forces
$p_{k}=-(1/L)g_{k}$, so the trust‑region step coincides with the standard
gradient step. The bound (7) reduces to the classic $L R^{2}/(2K)$ result.
\item \textbf{Non‑convex smooth case.}
If convexity is dropped but $f$ remains $L$‑smooth, the same analysis
yields a bound on the minimum gradient norm over the first $K$
iterations: $\min_{0\le k<K}\norm{\nabla f(x_{k})}^{2}=O(1/K)$.
\end{enumerate}

\end{document}